\documentclass[10pt]{article}
\usepackage[utf8]{inputenc}
\usepackage[T1]{fontenc}
\usepackage{amsmath}
\usepackage{amsfonts}
\usepackage{amssymb}
\usepackage[version=4]{mhchem}
\usepackage{stmaryrd}
\usepackage{graphicx}
\usepackage[export]{adjustbox}
\graphicspath{ {./images/} }

\begin{document}
Equilibrium refinements in extensive-form games

\section*{Subgame perfection and non-credible threats}
\begin{center}
\begin{tabular}{lll}
 & $\mathbf{T}$ & $\mathbf{B}$ \\
\hline
u & 2,2 & $[2,2]$ \\
\hline
d & $[3,1]$ & 1,0 \\
\hline
\end{tabular}
\end{center}

\begin{itemize}
  \item Two Nash equilibria
  \item $(d, T)$.
  \item $(u, B)$. Credible?
  \item Subgame perfection eliminates noncredible threats in games of perfect information.
\end{itemize}

\section*{Subgame of a game}
\begin{itemize}
  \item Let $x$ be a node in ( $X, \succ$ ), the game tree in an extensive form game $\Gamma$. The set of nodes $X^{\prime}=\{y: y \succ x\}$ is a subgame, denoted by $\Gamma(x)$, if $X^{\prime}$ with the structure inherited from $\Gamma$, is an extensive form game.
\end{itemize}

\section*{Remarks}
\begin{enumerate}
  \item The information set containing $x$ must be a singleton.
  \item For every information set $u$ in $\Gamma$, either $u \subset S$ or $u \cap X^{\prime}=\emptyset$.
\end{enumerate}

Let $\Gamma$ be an extensive form game with $I$ players.

\begin{itemize}
  \item For any node $x$, a behavior-stragtegy profile $\sigma=\left(\sigma_{1}, \ldots, \sigma_{I}\right)$ determines the probability $P_{\sigma}(x)$ that $x$ will be reached.
\end{itemize}

See Maschler, Solan, and Zamir (2013), page 254.\\
Definition 1 (Equilibrium path) Given $\Gamma$ and a behavior strategy profile $\sigma$, the path of play according to $\sigma$ is the set of nodes that are reached with strictly positive probability when players move according to $\sigma$.

If $\sigma$ is an equilibrium, the path of play is termed the equilibrium path.

\section*{Path of play and equilibrium path}
\section*{Matching pennies with outside option}
\begin{center}
\includegraphics[max width=\textwidth, alt={}]{4c633bb9-8b10-4905-874e-56542c304c8f-07_1202_1205_446_172}
\end{center}

Matching pennies. NE: $\sigma_{1}=(1 / 2,1 / 2)=\sigma_{2}$

\begin{itemize}
  \item Given $\Gamma, \mathcal{U}_{i}$ is the set of player $i$ 's information sets in $\Gamma$.
\end{itemize}

Definition 2 An information set $h_{i} \in \mathcal{U}_{i}$ is said to be in the path of play for a profile $\sigma$, if it is reached with positive probability according to $\sigma$, i.e.\\
$P_{\sigma}\left(h_{i}\right)>0$.

\section*{NE and Subgame Perfect equilibrium}
Theorem 1 If a strategy profile $\sigma$ is a NE of an extensive form game $\Gamma$, then $\sigma$ induces a NE in any subgame $\Gamma(x)$ where $x$ is in the equilibrium path of $\Gamma$ (that is to say $P_{\sigma}(x)>0$ ).

For instance, Maschler, Solan, and Zamir (2013), Theorem 7.5, page 254.\\
Definition 3 (Subgame perfect) A strategy profile $\sigma$ is a Subgame Perfect Equilibrium (SPE) of an extensive form game $\Gamma$ if $\sigma$ induces a NE in every subgame of $\Gamma$.

Theorem 2 Every SPE of $\Gamma$ is also a NE.

\section*{Sequential games and subgame perfection}
$$
\begin{array}{lll} 
& L_{2} & H_{2} \\
\hline L_{1} & 4,4 & 2,8 \\
\hline H_{1} & 8,2 & 1,1
\end{array}
$$

\includegraphics[max width=\textwidth, alt={}, center]{4c633bb9-8b10-4905-874e-56542c304c8f-10_1030_784_862_184}\\
\includegraphics[max width=\textwidth, alt={}, center]{4c633bb9-8b10-4905-874e-56542c304c8f-10_1042_787_850_1413}

Each player $i$ has a single information set; thus $S_{i}=\left\{L_{i}, H_{i}\right\}$.\\
\includegraphics[max width=\textwidth, alt={}, center]{4c633bb9-8b10-4905-874e-56542c304c8f-12_1055_797_193_177}\\
\includegraphics[max width=\textwidth, alt={}, center]{4c633bb9-8b10-4905-874e-56542c304c8f-12_1075_817_177_1405}

\section*{Subgame perfection-Example}
\begin{center}
\begin{tabular}{ccccc}
 & $\left(H_{2}, H_{2}\right)$ & $\left(H_{2}, L_{2}\right)$ & $\left(L_{2}, H_{2}\right)$ & $\left(L_{2}, L_{2}\right)$ \\
\hline
$H_{1}$ & 4,4 & 4,4 & 2,8 & $[2,8]$ \\
\hline
$L_{1}$ & $[8,2]$ & 1,1 & $[8,2]$ & 1,1 \\
\hline
\end{tabular}
\end{center}

\begin{center}
\includegraphics[max width=\textwidth, alt={}]{4c633bb9-8b10-4905-874e-56542c304c8f-13_1055_1479_844_218}
\end{center}

Strategy player 2: $\left(H_{2}, L_{2}\right)$ indicates $H_{2}$ if one played $H_{1}$.

\section*{Completely mixed}
\section*{Definition 4 (Completely mixed)}
\begin{itemize}
  \item A mixed strategy $\mu_{i}$ for player $i$ in a finite normal form game is completely mixed if $\mu_{i}(s)>0$ for every pure strategy $s$ of player $i$.
  \item A strategy profile $\mu$ is completely mixed if $\mu_{i}$ is completely mixed for every $i$.
  \item A behavioral strategy $\sigma_{i}$ for player $i$ in a finite extensive form game is completely mixed if for every information set $h \in \mathcal{U}_{i}$, and every action $a \in C(h), \sigma_{i}(a \mid h)>0$.
  \item A strategy profile $\sigma$ is completely mixed if $\sigma_{i}$ is completely mixed for every $i$.
\end{itemize}

\section*{Completely mixed NE and SPE}
Theorem 3 Every NE in completely mixed strategies of a finite extensiveform game is a SPE.

See Maschler, Solan, and Zamir (2013), Corollary 7.7.\\
Theorem 4 Every finite extensive-form game with perfect information has a SPE in pure strategies.

Proof is a backward induction argument.\\
See Maschler, Solan, and Zamir (2013), Theorem 7.9.

\section*{Example: Entry deterrence}
\includegraphics[max width=\textwidth, alt={}, center]{4c633bb9-8b10-4905-874e-56542c304c8f-17_1350_2401_409_178}\\
\includegraphics[max width=\textwidth, alt={}, center]{4c633bb9-8b10-4905-874e-56542c304c8f-18_1350_2398_261_181}

\section*{Grab the Dollar}
\begin{center}
\includegraphics[max width=\textwidth, alt={}]{4c633bb9-8b10-4905-874e-56542c304c8f-19_577_2521_566_227}
\end{center}

\section*{References}
Maschler, Michael, Eilon Solan, and Shmuel Zamir. 2013. Game Theory. Translated by Ziv Hellman. 1st ed. Cambridge University Press.


\end{document}